\documentclass[12pt,a4paper]{article}

\usepackage[utf8]{inputenc}
\usepackage[T1]{fontenc}
\usepackage[american]{babel}
\usepackage{amsmath,amssymb,amsthm}
\usepackage{booktabs}
\usepackage{graphicx}
\usepackage{algorithm}
\usepackage{algpseudocode}
\usepackage{hyperref}
\usepackage[margin=2.5cm]{geometry}

\title{Integer-Only 3D Rotation via Paeth Decomposition\\with Bresenham Accumulators}
\author{A.~F.~Neto}
\date{\today}

\begin{document}
\maketitle

\begin{abstract}
We present a cellular automaton (CA) capable of rotating the contents
of a three-dimensional cubic lattice of side~$L$ by an arbitrary angle
$\theta$ around the $z$-axis using exclusively integer arithmetic:
addition, subtraction, comparison, and bit-shift.  The angle is
specified by a pair of integers $(p,q)$ satisfying
$\sin\theta = p/q$.  The method decomposes the rotation into three
shear operations (Paeth decomposition), each executed by a Bresenham
accumulator that propagates from the grid center outward---replacing
all multiplications with cumulative additions.  We prove that (i)~the
mapping is bijective, preserving the number of active cells exactly,
and (ii)~the residual shear artifact vanishes as $O(1/L)$ when the
pair $(p,q)$ is scaled by~$L$, so the result converges to the ideal
rotation in the continuum limit.  Numerical experiments for
$L\in\{128,256,512\}$ confirm both properties.
\end{abstract}

% ==================================================================
\section{Introduction}
\label{sec:intro}

Discrete rotations on integer lattices arise naturally in cellular
automaton models of physics, where all state variables must be
representable as integers and the update rule may use only
addition, subtraction, comparison, and bit-shift---excluding
multiplication and division at the cell level~\cite{toffoli1987}.

The standard approach to integer-grid rotation is the
\emph{three-shear decomposition} introduced by
Paeth~\cite{paeth1986}: every 2D rotation can be written as the
composition of three axis-aligned shears, each of which is a
bijection on integer coordinates.  We extend this to 3D (rotation
around a fixed axis) and replace the per-row multiplication with a
Bresenham-style accumulator, yielding a procedure that is
compatible with CA constraints.

The main contributions are:
\begin{enumerate}
\item A complete integer-only rotation algorithm parameterized by
      a single pair $(p,q)$, with no floating-point arithmetic
      anywhere in the pipeline.
\item An $L$-scaling technique that makes the angular precision of
      the \texttt{isqrt} (integer square root) co-scale with the
      grid resolution, guaranteeing convergence to the exact
      rotation.
\item Numerical verification of bit-conservation and convergence
      for grids up to $L=512$ ($\approx1.3\times10^8$ cells).
\end{enumerate}

% ==================================================================
\section{Mathematical Framework}
\label{sec:math}

\subsection{Paeth three-shear decomposition}

A counter-clockwise rotation by angle $\theta$ in the $(x,y)$ plane
can be factored as
\begin{equation}
\label{eq:paeth}
R(\theta)
= \underbrace{\begin{pmatrix}1 & -t\\0 & 1\end{pmatrix}}_{S_x(-t)}
  \underbrace{\begin{pmatrix}1 & 0\\s & 1\end{pmatrix}}_{S_y(s)}
  \underbrace{\begin{pmatrix}1 & -t\\0 & 1\end{pmatrix}}_{S_x(-t)},
\end{equation}
where $s = \sin\theta$ and $t = \tan(\theta/2)$.  Since the
rotation is around~$z$, each $z$-slice is rotated independently;
the third coordinate is inert.

\subsection{Integer parameterization}

The angle is specified by the pair $(p,q)$ with
$\sin\theta = p/q$, \ $0 < p \le q$.  From this we derive:
\begin{align}
\cos\theta &\approx \frac{\mathrm{isqrt}(q^2 - p^2)}{q},
\label{eq:cos}\\[4pt]
\tan\!\tfrac{\theta}{2}
  &= \frac{\sin\theta}{1+\cos\theta}
   = \frac{p}{q + \mathrm{isqrt}(q^2 - p^2)}
   \;\equiv\; \frac{t_p}{t_q},
\label{eq:tan}
\end{align}
where $\mathrm{isqrt}(n) = \lfloor\sqrt{n}\rfloor$ is computed by
the classical bit-by-bit algorithm using only addition, subtraction,
comparison, and right-shift.

\subsection{$L$-scaling for precision}
\label{sec:scaling}

For small $(p,q)$ the truncation $\mathrm{isqrt}(q^2-p^2)$ can be
far from $\sqrt{q^2-p^2}$.  For example, $(p,q)=(1,2)$ gives
$\mathrm{isqrt}(3)=1$ versus $\sqrt{3}\approx1.732$, introducing a
24\% error in $\tan(\theta/2)$.

We eliminate this by scaling: before computing~\eqref{eq:tan}, replace
$(p,q)$ with $(pL,\,qL)$.  The sine ratio is unchanged
($pL/qL = p/q$), but the isqrt argument becomes $L^2(q^2-p^2)$,
whose integer square root has absolute error $\le 1$ and
\emph{relative} error $\le 1/(L\sqrt{q^2-p^2})$.  Thus
\begin{equation}
\label{eq:tanerr}
\left|\frac{t_p}{t_q} - \tan\!\tfrac{\theta}{2}\right|
= O\!\left(\frac{1}{L}\right).
\end{equation}

% ==================================================================
\section{Algorithm}
\label{sec:algo}

\subsection{Bresenham shift computation}

Each shear requires shifting row~$y$ (or column~$x$) by
$\lfloor d\cdot t_p/t_q \rfloor$ cells, where $d=|y - L/2|$ is the
distance from the grid center.  Rather than multiplying, we propagate
a Bresenham accumulator from $d=0$ outward:

\begin{algorithm}[H]
\caption{Compute cumulative shifts via Bresenham}
\label{alg:bres}
\begin{algorithmic}[1]
\State $\mathrm{acc} \gets 0$; \; $\mathrm{total} \gets 0$;
  \; $\mathrm{shifts}[\mathrm{center}] \gets 0$
\For{$d = 1, 2, \ldots, L-1$}
  \State $\mathrm{acc} \gets \mathrm{acc} + t_p$
  \While{$\mathrm{acc} \ge t_q$}
    \State $\mathrm{acc} \gets \mathrm{acc} - t_q$
    \State $\mathrm{total} \gets \mathrm{total} + 1$
  \EndWhile
  \State $\mathrm{shifts}[\mathrm{center}+d] \gets
         +\sigma\cdot\mathrm{total}$
  \State $\mathrm{shifts}[\mathrm{center}-d] \gets
         -\sigma\cdot\mathrm{total}$
\EndFor
\end{algorithmic}
\end{algorithm}

\noindent
Here $\sigma = \pm1$ encodes the shear direction.  The only
operations are addition, subtraction, and comparison.  The loop body
maps directly to a wavefront propagation rule in a CA: cell at
distance~$d$ reads its neighbor at distance~$d-1$ and applies one
Bresenham step.

\subsection{Shear execution}

Given the shift array, each shear is a bijective permutation within
each row (for $x$-shears) or column (for $y$-shears):
\begin{equation}
\mathrm{grid}_{\mathrm{new}}[x, y, z]
= \mathrm{grid}_{\mathrm{old}}[x - \mathrm{shifts}[y],\; y,\; z].
\end{equation}
Cells whose source index falls outside $[0,L)$ remain empty.

\subsection{Complete rotation}

The full rotation consists of three sequential shears:
\begin{enumerate}
\item $x$-shear with slope $-t_p/t_q$ along $y$,
\item $y$-shear with slope $+pL/(qL)$ along $x$,
\item $x$-shear with slope $-t_p/t_q$ along $y$ (identical to step~1).
\end{enumerate}

\subsection{Cellular automaton interpretation}

Each shear decomposes into two CA phases:
\begin{description}
\item[Propagation phase ($L/2$ ticks):]
  A Bresenham wavefront propagates from the center row (or column)
  outward, one cell per tick.  Each cell reads its inward
  neighbor's accumulator and applies a single add-compare step.
  This is structurally identical to the BFS wavefront in our
  pulsating-shell CA~\cite{neto2026arithmetico}.
\item[Shift phase ($\le L/2$ ticks):]
  All rows shift their content simultaneously, one cell per tick.
  Each cell with $\mathrm{shift\_remaining} > 0$ copies from its
  lateral neighbor and decrements.  Rows that finish early are inert.
\end{description}
Total ticks per rotation: $3\times(L/2 + L/2) = 3L$.

% ==================================================================
\section{Results}
\label{sec:results}

The algorithm is fully generic: it rotates \emph{any} configuration
of active cells, regardless of shape or topology.  For the
experiments below we use a centered cube of side~$L/2$ with a small
asymmetric flag (3~cells protruding in~$+x$) as a fiducial marker,
chosen solely because its sharp corners make geometric distortions
easy to detect.  The rotation angle is $\theta=30^\circ$, specified
by $(p,q)=(1,2)$.

\subsection{Bit conservation}

\begin{table}[h]
\centering
\caption{Active-cell counts before and after rotation.}
\label{tab:bits}
\begin{tabular}{@{}rrrc@{}}
\toprule
$L$ & Cells before & Cells after & Preserved \\
\midrule
128  & 262\,156 & 262\,156 & Yes \\
256  & 2\,097\,164 & 2\,097\,164 & Yes \\
512  & 16\,777\,228 & 16\,777\,228 & Yes \\
\bottomrule
\end{tabular}
\end{table}

The three-shear decomposition is a composition of bijective
permutations (integer translations within each row/column), so the
total number of active cells is preserved \emph{exactly} for any
angle, any~$L$, and any initial configuration---the result does not
depend on the shape of the object being rotated.

\subsection{Convergence of the tangent approximation}

Table~\ref{tab:tan} shows the error in $\tan(\theta/2)$ as a function
of the pair magnitude, both with and without $L$-scaling.

\begin{table}[h]
\centering
\caption{Tangent-half-angle precision.  The exact value is
$\tan15^\circ=0.26795$.}
\label{tab:tan}
\begin{tabular}{@{}rrrrr@{}}
\toprule
$p$ & $q$ & $\mathrm{isqrt}$ & $t_p/t_q$ & Error \\
\midrule
\multicolumn{5}{@{}l}{\emph{Without $L$-scaling (raw pair):}} \\
1 & 2 & 1 & 0.3333 & 24.4\% \\
5 & 10 & 8 & 0.2778 & 3.7\% \\
50 & 100 & 86 & 0.2688 & 0.32\% \\
500 & 1000 & 866 & 0.2680 & 0.001\% \\
\midrule
\multicolumn{5}{@{}l}{\emph{With $L$-scaling, $(p,q)=(1,2)$:}} \\
128 & 256 & 221 & 0.2683 & 0.15\% \\
256 & 512 & 443 & 0.2680 & 0.04\% \\
512 & 1024 & 886 & 0.2681 & 0.04\% \\
\bottomrule
\end{tabular}
\end{table}

\subsection{Geometric fidelity}

To quantify the residual shear, we de-rotate the output by~$-30^\circ$
(using floating-point, for measurement only) and measure the tilt
of the top edge and the maximum positional deviation from the ideal
square.

\begin{table}[h]
\centering
\caption{Geometric fidelity with $L$-scaled pair $(1,2)$.}
\label{tab:fidelity}
\begin{tabular}{@{}rrrr@{}}
\toprule
$L$ & Edge tilt & Max pos.\ error & Relative error \\
\midrule
128 & $0.29^\circ$ & 1.8 cells & 5.6\% \\
256 & $0.14^\circ$ & 1.6 cells & 2.5\% \\
512 & $0.07^\circ$ & 1.7 cells & 1.3\% \\
\bottomrule
\end{tabular}
\end{table}

The edge tilt and relative error both decrease as $O(1/L)$, confirming
convergence to the exact rotation.  The absolute positional error
remains bounded at $\approx 1.7$ cells (the inherent Bresenham
rounding of $\pm1$ cell per shear, compounded over three shears).

% ==================================================================
\section{Discussion}
\label{sec:discussion}

\subsection{Sources of error}

Two distinct sources contribute to the residual distortion:

\begin{enumerate}
\item \textbf{Tangent approximation error.}
  The integer square root in~\eqref{eq:cos} truncates
  $\sqrt{q^2-p^2}$ to $\lfloor\sqrt{q^2-p^2}\rfloor$, introducing
  a systematic bias in $\tan(\theta/2)$.  For $(p,q)=(1,2)$,
  $\mathrm{isqrt}(3)=1$ versus $\sqrt{3}=1.732$, yielding a 24\%
  error.  The $L$-scaling technique (Section~\ref{sec:scaling})
  reduces this to $O(1/L)$.

\item \textbf{Bresenham discretisation.}
  Each shear rounds the shift to the nearest integer, introducing
  $\pm 1$ cell error per row per shear.  After three shears, the
  maximum positional error is bounded by 3~cells regardless of~$L$.
  The \emph{relative} error is therefore $O(1/L)$.
\end{enumerate}

Both sources vanish in the limit $L\to\infty$, and the rotation
converges to the exact continuous result.

\subsection{CA compatibility}

The algorithm satisfies the constraints of our integer cellular
automaton framework:
\begin{itemize}
\item \textbf{No multiplication or division} in the cell update rule.
  The Bresenham accumulator replaces $d \times t_p / t_q$ with
  $d$~iterations of add-compare.  The $\mathrm{isqrt}$ uses only
  add, subtract, compare, and right-shift.
\item \textbf{Local neighborhood.}
  During the propagation phase, each cell reads one neighbor
  (toward the center) for the accumulator.  During the shift phase,
  each cell reads one lateral neighbor for the content.
\item \textbf{Deterministic and reversible.}
  Each shear is a bijection.  The inverse rotation uses the same
  three shears in reverse order with negated slopes.
\end{itemize}

\subsection{Angle specification}

The rotation angle is fully determined by the integer pair
$(p,q)$.  This is analogous to the Bresenham ratio
$(\mathrm{sinc\_p}, \mathrm{sinc\_q})$ in our sinc-wave CA: each
cell stores two integers that define a frequency (or, here, an
angle).  When all cells carry the same pair, the rotation is
uniform.  Non-uniform pairs would produce spatially varying
rotations---a mechanism for implementing curvature or torsion in
future work.

% ==================================================================
\section{Conclusion}
\label{sec:conclusion}

We have demonstrated that a three-dimensional rotation by an
arbitrary angle can be performed on an integer lattice using only
addition, subtraction, comparison, and bit-shift---operations
available to a cellular automaton.  The Paeth three-shear
decomposition, combined with Bresenham accumulators, provides an
exact bijection (zero bit loss) with geometric error that vanishes
as $O(1/L)$.  The $L$-scaling technique ensures that even the
simplest angle specification, such as $(1,2)$ for $30^\circ$, achieves
sub-degree accuracy for grids of modest size.

The algorithm is implemented in~\texttt{ac\_rotation.c}
(approximately 240 lines of~C) and is available in the project
repository.

% ==================================================================
\begin{thebibliography}{9}

\bibitem{paeth1986}
A.~W.\ Paeth,
``A fast algorithm for general raster rotation,''
in \emph{Graphics Interface '86}, 1986, pp.~77--81.

\bibitem{toffoli1987}
T.~Toffoli and N.~Margolus,
\emph{Cellular Automata Machines: A New Environment for Modeling},
MIT Press, 1987.

\bibitem{neto2026arithmetico}
A.~F.~Neto,
``Exact Euclidean squared distance reconstruction by local integer
propagation,''
\emph{ResearchGate preprint}, 2026.

\end{thebibliography}

\section*{Funding and competing interests}

The author declares that no funds, grants, or other support were received
during the preparation of this manuscript.
The author has no relevant financial or non-financial interests to disclose.

\end{document}
